\begin{titlepage}

    \begin{center}

        %% Extra whitespace at the top.
        \vspace*{2\bigskipamount}

        %% Print the title.
        {\makeatletter
            \titlestyle\bfseries\LARGE\@title
            \makeatother}

        %% Print the optional subtitle.
        {\makeatletter
            \ifx\@subtitle\undefined\else
                \bigskip
                \titlefont\titleshape\Large\@subtitle
            \fi
            \makeatother}

    \end{center}

    \cleardoublepage
    \thispagestyle{empty}

    \begin{center}

        %% The following lines repeat the previous page exactly.

        \vspace*{2\bigskipamount}

        %% Print the title.
        {\makeatletter
            \titlestyle\bfseries\LARGE\@title
            \makeatother}

        %% Print the optional subtitle.
        {\makeatletter
            \ifx\@subtitle\undefined\else
                \bigskip
                \titlefont\titleshape\Large\@subtitle
            \fi
            \makeatother}

        \vfill

        %% Apart from the names and dates, the following text is dictated by the
        %% promotieregelement.

        {\Large\titlefont\bfseries Proefschrift}

        \bigskip
        \bigskip

        ter verkrijging van de graad van doctor\\
        aan de Radboud Universiteit Nijmegen\\
        op gezag van de rector magnificus prof. dr. J.H.J.M. van Krieken,\\
        volgens besluit van het college voor promoties\\
        in het openbaar te verdedigen op\\

        \bigskip
        \bigskip

        door

        \bigskip
        \bigskip

        \makeatletter
        {\Large\titlefont\bfseries\@firstnames\ \MakeUppercase{\titleshape\@lastname}}
        \makeatother

        \bigskip
        \bigskip

        geboren op 28 mei, 1995\\
        te Tamil Nadu, India\\

        %% Extra whitespace at the bottom.
        \vspace*{2\bigskipamount}

    \end{center}

    \clearpage
    \thispagestyle{empty}

    \noindent This dissertation has been approved by the promotors.

    \bigskip
    \noindent Composition of the doctoral committee:
    
    \begin{tabbing}
        \hspace{\tabcolsep}\=\hspace{0.33\textwidth}\=\hspace{0.66\textwidth}                   \\[-3\medskipamount]
        \> Prof.\ dr.\ Britta \ Redlich,    \> Radboud Universiteit, \textit{promotor}      \\
        \> Dr.\ Sandra \ Br\"unken,        \> Radboud Universiteit, \textit{copromotor}    \\[\medskipamount]
        \>\textit{Manuscript committee:}  \\[\smallskipamount]
        \> dr. \ Oskar Asvany            \> Universit\"at zu K\"oln \\
        \> prof. \ dr.\ Jana \ Roithova            \> Radboud Universiteit \\
        \> prof. \ dr.\ Melanie \ Schnell   \> Universit\"at zu Kiel and DESY \\[\medskipamount]
        \>\textit{Other members:} \\[\smallskipamount]
    \end{tabbing}

    \medskip

    \vfill
    \begin{center}
        \includegraphics[height=0.5in]{_logos/Logo_RU_NL_RGB.pdf}
        \hspace{2em}
        \includegraphics[height=0.5in]{_logos/NWO.jpg}
    \end{center}
    \vfill
    \noindent

    \begin{tabular}{@{}p{0.2\textwidth}@{}p{0.8\textwidth}@{}}
        \textit{Keywords:}    &  astrochemistry, molecular ions, cryogenic ion trap, vibrational spectra, rotational spectra, Renner-Tellar effects \\[\medskipamount]
        \textit{Printed by:}   & \\[\medskipamount]
        \textit{Cover by:} & {
            \makeatletter
            \@initials~\@lastname
            \makeatother
        }
    \end{tabular}

    \vspace{4\bigskipamount}

    \noindent Copyright \textcopyright{} \the\year{} by{
        \makeatletter
        \@initials~\@lastname
        \makeatother
    }

    \medskip
    \noindent ISBN 000-00-0000-000-0

    \medskip
    % \noindent An electronic copy of this dissertation is available at\\
    % \url{https://repository.ru.nl/ ?}.
    This work is part of the research programme "ROSAA" with project number 740.018.010 and “HFML-FELIX: a 
    Dutch Centre of Excellence for Science under Extreme Conditions” (with project number 184.035.011) of the research 
    programme “Nationale Roadmap Grootschalige Wetenschappelijke Infastructuur”, which is (partly) financed by the 
    Netherlands Organisation for Scientific Research (NWO).
\end{titlepage}

